\subsection{Metodología Extreme Programming (XP)}
Es una metodología ágil que permite a los desarrolladores responder con confianza a los requisitos cambiantes de los clientes, promueve el trabajo en equipo, se basa en una realimentación continua entre el cliente y el equipo de desarrollo y una comunicación fluida entre los participantes \cite{Canos2003}.

\subsubsection{Ciclo de Vida de Extreme Programming} 
El ciclo de vida ideal de XP consiste en las siguientes fases \cite{Joskowicz2014}:
\begin{itemize}
    \item \textbf{Fase I: Exploración} \ \\
Se realizan las historias de usuarios, al mismo tiempo el equipo de desarrollo se adecua a las herramientas, tecnologías y prácticas que se implementarán para el desarrollo del proyecto. 

 \item\textbf{Fase II: Planificación de la Entrega } \ \\
Se establece la prioridad de cada historia de usuario y estima el esfuerzo para cada una de ellas, y se determina un cronograma en conjunto con el cliente, el resultado de esta fase es un Plan de Entregas, o “Release Plan”.

 \item\textbf{Fase III: Iteraciones } \ \\
Es la fase donde se desarrollan todas las funcionalidades, la salida es un entregable funcional que se le presenta al cliente, cada iteración dura de 1 a 4 semanas.

 \item\textbf{Fase IV: Producción} \ \\
Es la fase donde se realiza las pruebas adicionales y revisiones de rendimiento, al mismo tiempo se tiene que tomar decisiones sobre la incorporación de nuevas características al sistema antes de que el sistema sea entregado al cliente.

\end{itemize}\ \\
En cada iteración se realiza un ciclo completo de análisis, diseño, desarrollo y pruebas, utilizando un conjunto de reglas y prácticas que caracterizan a XP. Entre algunas reglas y prácticas  que caracterizan a XP están:

\begin{itemize}

\item Historias de usuario: herramienta que sirve para dar a conocer los requerimientos del sistema.

\item Plan de entregas (“Release Plan”): es el cronograma del proyecto, agrupa historias de usuario para conformar entregas.

\item Plan de Iteraciones (Iteration Plan): las historias agrupadas en el “release plan” se van desarrollando en iteraciones. 

\item Reuniones Diarias de Seguimiento (Stand – Up Meeting), mantener una constante comunicación con el equipo de desarrollo para así comunicar problemas y soluciones.

\item Simplicidad: un diseño sencillo siempre es mejor que uno complejo. 

\item Recodificación: escribir de nuevo el código de un programa que tenga la misma funcionalidad, con el objetivo de que pueda ser más simple o entendible.

\item Programación dirigida por las pruebas: Antes de codificar se deben primero escribir las pruebas que el sistema debe pasar.

\item Programación en pares: Dos programadores juntos en un mismo computador para para minimizar errores y lograr mejores diseños.

\item Detección y corrección de errores: si se encuentra un error, se debe corregir inmediatamente.

\item  Pruebas de aceptación: cada historia de usuario debe tener una prueba de aceptación y se debe de realizar después de cada iteración con el objetivo de comprobar si la historia de usuario fue implementada correctamente.
\end{itemize}

\subsubsection{Roles en Extreme Programming} 
Los Roles que se han encontrado efectivos en la Programación Extrema son \cite{Canos2003}: 
\begin{itemize}
\item\textbf{Desarrollador: } El rol del desarrollador es el más importante en Extreme Programming, tiene como responsabilidad: estimar historias, definir tareas de historias, estimar tareas, escribir pruebas de aceptación, además, debe tener como habilidades: saber comunicarse y codificar solo lo que requiere. El desarrollador desempeña distintas funciones como: programador, arquitecto y diseñador, arquitecto de interfaz / diseñador de interfaz de usuario, diseñador de bases de datos y probador. 

\item\textbf{Cliente:} El rol del cliente es tan crucial como el rol del desarrollador, es el que debe saber qué programar, por lo tanto, debe tener ciertas habilidades como: tener una actitud positiva para el éxito del proyecto, influir en el proyecto sin poder controlarlo, tomar decisiones que se requieren de vez en cuando, cambiar las decisiones a medida que el producto evoluciona. También es necesario, que el cliente esté en comunicación constante con el equipo.

\item\textbf{Tracker:} Es el encargado de supervisar el progreso del desarrollo del software. Por lo general, revisa el progreso de los programadores y ofrece ayuda donde sea necesario de varias maneras.

\item\textbf{Manager:} Es el encargado de planificar las reuniones, asiste a las reuniones y registra los resultados de ellas, mantiene al equipo productivo. 

\end{itemize}